We want to simplify a mesh while preserving its \emph{mechanical responsiveness},
formulating a quadric error analogous to the geometric QEM of Garland \& Heckbert.
This section builds the analogy from the ground up, identifies the natural
``plane quadric'' for the mechanical case, and is honest about where the geometric
and mechanical settings genuinely diverge.

\subsection{What the QEM plane quadric actually is}

Stripped to its skeleton, the geometric face quadric is a \textbf{squared linear residual}:
\begin{equation}
  E_f(\mathbf{x}) = r_f(\mathbf{x})^2, \qquad r_f(\mathbf{x}) = \mathbf{n}^T\mathbf{x} + d.
\end{equation}
\begin{itemize}
  \item The \textbf{variable} is an absolute point position $\mathbf{x}\in\mathbb{R}^3$.
  \item The \textbf{residual} $r_f$ is the signed distance to the plane---a \emph{linear
        functional} of the variable that is zero on the ``correct'' set (the plane itself).
  \item The matrix $\mathbf{n}\mathbf{n}^T$ is rank-1 and positive semi-definite; its
        \textbf{null space is the in-plane directions}---the motions that cost nothing.
  \item Accumulation ($\sum_f$) works because every face's residual is written in the
        \emph{same} variable $\mathbf{x}$.
\end{itemize}
So QEM follows a template: \emph{pick a quantity that should be preserved (flatness),
write its violation as a squared linear residual, sum the residuals.}

\subsection{Transplanting the template to mechanics}

The quantity we want to preserve is \emph{elastic response}, encoded by the element
energy. For the CST membrane element,
\begin{equation}
  W_e(\mathbf{d}) = \tfrac12\,\mathbf{d}^T K_e\,\mathbf{d}
                  = \tfrac12\,A\,t\,\boldsymbol\varepsilon^T D\,\boldsymbol\varepsilon,
  \qquad \boldsymbol\varepsilon = B\mathbf{d}.
\end{equation}
This is \emph{already} a squared linear residual, but with a non-trivial metric.
Factor the constitutive matrix as $D = L L^T$ (Cholesky, i.e.\ $D^{1/2}$) and define
$M = \sqrt{A t}\,L^T B \in \mathbb{R}^{3\times 9}$. Then
\begin{equation}
  W_e(\mathbf{d}) = \tfrac12\sum_{k=1}^{3}\big(\mathbf{m}_k^T\mathbf{d}\big)^2,
\end{equation}
where $\mathbf{m}_k^T$ are the rows of $M$. \textbf{Each term is a generalized plane
quadric.} The strain plays the role of distance, and the constitutive matrix $D$ is the
metric that geometric QEM leaves as the identity.

The bending hinge is even cleaner---it is \emph{literally} a single plane quadric:
\begin{equation}
  W_b(\mathbf{d}) = \tfrac12\,c\,(\mathbf{g}^T\mathbf{d})^2
                  = \tfrac12\,\big(\underbrace{\sqrt{c}\,\mathbf{g}}_{\text{``normal''}}{}^{\!T}\mathbf{d}\big)^2.
\end{equation}
The linearized dihedral angle $\mathbf{g}^T\mathbf{d}$ \emph{is} the signed-distance
residual; $\sqrt{c}\,\mathbf{g}$ \emph{is} the plane normal---only it lives in $12$-D
displacement space instead of $3$-D position space.

\subsection{The dictionary}

\begin{center}
\renewcommand{\arraystretch}{1.3}
\begin{tabular}{|p{0.46\linewidth}|p{0.46\linewidth}|}
\hline
\textbf{QEM (geometry)} & \textbf{Mechanical analog} \\
\hline
variable: point position $\mathbf{x}\in\mathbb{R}^3$
  & nodal displacement $\mathbf{d}\in\mathbb{R}^{9}$ (or $\mathbb{R}^{12}$ hinge) \\
\hline
residual: signed distance $\mathbf{n}^T\mathbf{x}+d$
  & strain $\boldsymbol\varepsilon = B\mathbf{d}$ (resp.\ $\mathbf{g}^T\mathbf{d}$) \\
\hline
metric on residual: identity
  & constitutive $D$ (scaled by $A t$) \\
\hline
face quadric: $\mathbf{n}\mathbf{n}^T$, rank $1$
  & element stiffness $K_e = A t\,B^T D B$, rank $3$ (hinge: rank $1$) \\
\hline
null space: in-plane slides (don't matter)
  & rigid-body modes (don't matter) \\
\hline
accumulation: $\sum_f$ on shared $\mathbf{x}$
  & assembly $K = \sum_e K_e$ on shared DOFs \\
\hline
\end{tabular}
\end{center}

\noindent\textbf{The mechanical ``plane quadric'' is the element stiffness $K_e$}, viewed
as a squared strain-residual in the energy metric $D$. It is positive semi-definite, its
null space is exactly the deformations that ``don't matter'' (rigid motion), and it
factors into rank-1 generalized plane quadrics whose ``normals'' are strain-sensing
co-modes. This is not a hack---it is the same template (preserve-a-quantity
$\rightarrow$ squared-linear-residual $\rightarrow$ sum) instantiated on strain instead
of distance.

\subsection{Where it genuinely diverges}

Two differences are not cosmetic; they are where the real design work lives.

\paragraph{1. The variable is a displacement (relative), not a position (absolute).}
Geometric fidelity is a \emph{zeroth-order} property of one point; mechanical fidelity is
a \emph{first-order} property---strain is a gradient, a \emph{difference} of
displacements. A single vertex's displacement determines no strain by itself. This is why
geometric QEM can collapse a vertex's incident faces into one portable $3\times3$ form on
that vertex's coordinate, but mechanics cannot: each element's quadric lives on its whole
stencil, and adjacent quadrics couple through shared DOFs. The additive structure still
exists---it is finite-element assembly---but it lives on the \emph{full} displacement
vector and never reduces to a per-vertex form. That coupling \emph{is} the mechanics; any
reduction that discards it discards what we are trying to preserve.

\paragraph{2. The face quadric is a \emph{response template}, not yet an \emph{error}.}
In geometric QEM the same quadric does double duty: it describes the plane \emph{and},
evaluated at the moved point, yields the error---because the original surface has zero
distance to its own planes, so any post-collapse distance \emph{is} the error.
Mechanically, $K_e$ describes the response, but the simplification error is the
\textbf{discrepancy} between fine and coarse response, which is not zero-referenced.
Concretely, with a prolongation $S$ mapping coarse displacements to fine ones, the
mechanically-exact coarse operator is the Galerkin form
\begin{equation}
  K_c = S^T K_f\,S,
\end{equation}
and the collapse error is how far the actual coarse $K_c$ deviates from it---measured
\emph{on some set of deformations}.

That last phrase is the real fork: \textbf{which deformations do we probe the discrepancy
on?} This is the mechanical counterpart of ``which planes matter,'' and it is where
``responsiveness'' gets \emph{defined}:
\begin{itemize}
  \item \textbf{affine / linear displacement modes} $\rightarrow$ preserves the effective
        (homogenized) stiffness;
  \item \textbf{low-frequency eigenmodes of $K$} $\rightarrow$ preserves global compliance
        / vibration response;
  \item \textbf{a known load subspace} $\rightarrow$ preserves response to the loads we
        actually care about.
\end{itemize}
Choosing that probe subspace is what eventually lets us reduce the coupled assembly back
down to a small, portable, QEM-shaped quadric---i.e.\ it is the principled route to
additivity, rather than a hack imposed up front.
The remaining subsections develop this
program: we first establish \emph{why} the relative-displacement quadric resists
per-vertex reduction, then identify the two tools that buy reducibility back---static
condensation (Schur complement) for exact static fidelity, and a probe subspace for true
QEM portability---and finally show how the probe choice dials the mechanical quadric
between geometric QEM and an attribute-quadric.

\subsection{Why the mechanical quadric resists per-vertex reduction}

Geometric QEM reduces to a portable $4\times4$ form per vertex for two reasons, and
mechanics violates both.
\begin{itemize}
  \item \textbf{The residual is linear in a \emph{single} vertex's coordinate.}
        $r_f(\mathbf{x}) = \mathbf{n}^T\mathbf{x}+d$ involves one point. Strain
        $\boldsymbol\varepsilon = B\mathbf{d}$ is a discrete gradient---a \emph{difference}
        $\sim(\mathbf{u}_j-\mathbf{u}_i)/\ell$. A single node's displacement determines no
        strain, so the smallest object on which the residual is even defined is the element
        stencil, not a vertex.
  \item \textbf{The coefficients are frozen.} The plane $(\mathbf{n},d)$ is a constant
        reference baked in once at preprocessing. Mechanics has no frozen target
        displacement: the rest state stores zero energy, but so does \emph{every} rigid
        motion of it. The ``correct set'' is a whole null space (the rigid-body modes), not
        a point.
\end{itemize}
So the honest answer is that the mechanical quadric \textbf{does not reduce to a per-vertex
form in general}---this is exactly divergence~1 above. The quadric is irreducibly the
element form $K_e\in\mathbb{R}^{9\times9}$, and ``accumulation'' is finite-element assembly
$K=\sum_e K_e$ on the \emph{global} DOF vector. Reducibility can only be bought back in one
of two ways:
\begin{itemize}
  \item \textbf{Eliminate the vanishing DOF exactly} (Schur complement, \S\,below): keeps
        the full static response on the survivors; pays in fill-in.
  \item \textbf{Restrict to a frozen, per-vertex-factorable reference} (probe subspace,
        \S\,below): recovers a literal per-vertex additive quadric; pays in off-probe
        response.
\end{itemize}

\subsection{Schur complement: the mechanically-exact collapse operator}

An edge collapse \emph{removes} degrees of freedom, and static condensation is the exact
way to remove DOFs while preserving the response on what remains. Partition the global DOFs
into \emph{retained} ($r$) and \emph{eliminated} ($e$, the disappearing vertex's three
DOFs),
\begin{equation}
  K = \begin{pmatrix} K_{rr} & K_{re} \\ K_{er} & K_{ee} \end{pmatrix}.
\end{equation}
With no load on the eliminated DOFs they equilibrate to
$\mathbf{d}_e = -K_{ee}^{-1}K_{er}\mathbf{d}_r$, and the energy becomes
\begin{equation}
  \tfrac12\mathbf{d}^T K\mathbf{d}
    = \tfrac12\,\mathbf{d}_r^T\,\widehat{K}_{rr}\,\mathbf{d}_r,
  \qquad
  \widehat{K}_{rr} = K_{rr} - K_{re}K_{ee}^{-1}K_{er},
\end{equation}
where $\widehat{K}_{rr}$ is the \textbf{Schur complement}. Three facts make this the right
object:
\begin{itemize}
  \item \textbf{It is the Galerkin operator with the \emph{optimal} prolongation.}
        $\widehat{K}_{rr} = S^T K_f S$ when $S$ is the \emph{harmonic extension}
        (eliminated DOFs interpolated to minimize energy). So the Schur complement is not an
        alternative to $K_c = S^T K_f S$; it is the best-case $K_c$. Any actual collapse
        uses a simpler prolongation, and the collapse error is precisely
        $\|S_{\text{actual}}^T K_f S_{\text{actual}} - \widehat{K}_{rr}\|$.
  \item \textbf{It is local.} $K_{er}$ is nonzero only on the $1$-ring (FEM sparsity), so
        $\widehat{K}_{rr}$ differs from $K_{rr}$ only inside the $1$-ring. Eliminating a
        vertex is a $1$-ring operation---the locality QEM requires.
  \item \textbf{The error has a name: fill-in.} The correction
        $-K_{re}K_{ee}^{-1}K_{er}$ is a dense block coupling the neighbors of the
        eliminated vertex---Gaussian-elimination fill-in. Structurally, ``mechanical QEM''
        is approximate sparse Cholesky / algebraic-multigrid coarsening, and the quadric
        error of a collapse is the magnitude of the fill-in one refuses to store.
\end{itemize}

\subsection{Edge collapse as common-mode retention}

The collapse $(v_i,v_j)\to v^\star$ leaves the \emph{link} (the boundary of the edge's
$1$-ring) unchanged: same outer vertices, same outer edges. Only the interior is re-wired.
Mechanically, the boundary DOFs are fixed shared data, so the collapse is a local problem
with the link as a \textbf{Dirichlet boundary} and the merged vertex as the only free
interior DOF---exactly what the Schur complement needs.

To make ``the two vertices wiggle as one'' precise, change basis on the pair into a
\textbf{common mode} and a \textbf{differential mode},
\begin{align*}
  \mathbf{d}_+ &= \tfrac12(\mathbf{d}_i + \mathbf{d}_j)
    \quad(\text{move together}),
  \\
  \mathbf{d}_- &= \mathbf{d}_i - \mathbf{d}_j
    \quad(\text{stretch/shear the edge}).
\end{align*}
Two free vertices ($6$ DOF) split into $3$ common $+$ $3$ differential. The collapse
\textbf{keeps the common mode} ($\mathbf{d}^\star\leftrightarrow\mathbf{d}_+$) and
\textbf{removes the differential mode}---this is the $6\to3$ reduction. Hence the faithful
condition:
\begin{quote}
  The collapse is faithful iff, on the probe deformations, the differential mode
  $\mathbf{d}_-$ is unexcited---i.e.\ $v_i$ and $v_j$ were already moving together. The
  collapse cost is the elastic energy that lived in $\mathbf{d}_-$.
\end{quote}
Under a smooth probe field, adjacent vertices displace almost equally, so
$\mathbf{d}_-\approx 0$ and the collapse is cheap; this is why short edges in smooth regions
coarsen safely. It is also exactly algebraic-multigrid aggregation: lump together DOFs whose
relative motion is unexcited.

There are two distinct ways to remove $\mathbf{d}_-$, and the gap between them \emph{is} the
error:
\begin{itemize}
  \item \textbf{Restriction (what the collapse literally does):} set $\mathbf{d}_-=0$ and
        duplicate, $\mathbf{d}_i=\mathbf{d}_j=\mathbf{d}^\star$. Freezing a DOF can only
        \emph{add} stiffness, so a collapsed edge is always at least as stiff as the
        original on any retained motion.
  \item \textbf{Condensation (the ideal):} let $\mathbf{d}_-$ relax to its
        energy-minimizing value given the boundary $\mathbf{d}_b$ and common mode,
        \begin{equation}
          \mathbf{d}_-^\star
            = -K_{--}^{-1}\big(K_{-+}\mathbf{d}_+ + K_{-b}\mathbf{d}_b\big),
        \end{equation}
        which preserves the static response exactly (the Schur complement of the
        differential block).
\end{itemize}
The error is the gap between $\mathbf{d}_-=0$ and $\mathbf{d}_-=\mathbf{d}_-^\star$---the
energy recoverable by relaxing rather than freezing the edge. The merged position
$\mathbf{x}^\star$ (and the prolongation) is the one knob: choose it so the single vertex's
wiggle reproduces what the pair actually did, i.e.\ so that $\mathbf{d}_-^\star$ is small on
the probe. That \emph{is} the QEM minimization, restated.

\paragraph{Caveat.} ``Boundary same, interior re-wired'' is right for connectivity, but
three energy terms change, and only the first is the differential mode: (i) the relative
stretch/shear stored across the ring (the $\mathbf{d}_-$ energy---the main term); (ii) the
two collapsing triangles' own membrane and hinge energy (they go to zero area and vanish);
(iii) the dihedral angles of every ring edge shift, so the bending references in the ring
move. The differential-mode picture captures (i); (ii)--(iii) are why bending wants to be
re-evaluated at the collapse step, not frozen.

\subsection{Stretching and bending: triangle quadrics versus edge quadrics}

The two energies separate cleanly by stencil:
\begin{center}
\renewcommand{\arraystretch}{1.3}
\begin{tabular}{|p{0.26\linewidth}|p{0.16\linewidth}|p{0.22\linewidth}|p{0.16\linewidth}|}
\hline
\textbf{energy} & \textbf{simplex} & \textbf{stencil} & \textbf{rank} \\
\hline
stretching (CST membrane) & triangle & 3 verts / 9 DOF & $3$ \\
\hline
bending (dihedral hinge) & interior edge & 4 verts / 12 DOF & $1$ \\
\hline
\end{tabular}
\end{center}
\noindent So bending enters at the \textbf{edge step}, as its own element family---not at
the triangle step. Two ways to handle it:
\begin{itemize}
  \item \textbf{Assembly view (recommended).} Assemble
        $K = \sum_{\text{tri}} K_e^{\text{stretch}} + \sum_{\text{edge}} K_e^{\text{bend}}$.
        Both families write into the $1$-ring blocks of the global operator, and the
        collapse (Schur or probe-reduction) operates on the \emph{merged} local operator.
        There is no ``triangle vs.\ edge step'' to decide because assembly already unifies
        them.
  \item \textbf{Garland-style per-vertex accumulation (if wanted for scoring).} Stretching
        distributes to its $3$ vertices like standard QEM; bending distributes to the
        hinge's $4$ stencil vertices (or its $2$ endpoints). But bending must be
        \emph{re-evaluated} at the collapse step, not frozen: a triangle's plane is
        collapse-invariant, whereas collapsing one edge perturbs the dihedral angle of every
        edge in the $1$-ring. Bending is therefore more fragile than stretching, with a
        wider stencil and a moving reference.
\end{itemize}
A deeper point connects to the probe discussion: \textbf{bending is invisible to
constant-strain (affine) deformation.} A flat patch under uniform in-plane strain does not
bend; bending is a second-derivative (curvature) quantity. So whether bending enters the
cost at all depends on the probe subspace---preserving bending requires enriching the probe
with quadratic (constant-curvature) modes $\mathbf{u}(X)=\tfrac12 X^T H X + FX + \mathbf{c}$.

\subsection{Affine probes and the homogenization patch test}

Take the probe subspace to be the \textbf{affine displacement fields}
$\mathbf{u}(X) = FX + \mathbf{c}$ with $F\in\mathbb{R}^{3\times3}$ constant. A constant-$F$
field produces \emph{constant strain} everywhere; the symmetric part of $F$ is strain, the
skew part is infinitesimal rotation, and $\mathbf{c}$ is translation (the last two are
rigid, zero-energy). The energetic content is the six independent constant-strain modes,
and the energy stored under a uniform macroscopic strain $\bar{\boldsymbol\varepsilon}$
\emph{is the definition} of the effective modulus,
\begin{equation}
  \tfrac12\,\bar{\boldsymbol\varepsilon}^T C^{\text{eff}}\,\bar{\boldsymbol\varepsilon}
    \cdot \mathrm{Vol} = \text{(total stored energy)}.
\end{equation}
Concretely, let the columns of $A$ be the affine modes sampled at the nodes. Then
\begin{equation}
  \text{``preserve effective stiffness''}
  \Longleftrightarrow
  A^T K_c A = A^T K_f A \ \text{on the patch,}
\end{equation}
which is precisely the finite-element \textbf{patch test}. This is the weakest but most
robust target: a six-dimensional subspace guaranteeing the coarse mesh is neither
artificially stiffer nor softer in bulk. It is blind to bending (see above), which is why
the eigenmode and load-subspace probes of the dictionary are needed when bending or
task-specific compliance matters.

\subsection{Reduction to an additive per-vertex quadric}

The payoff: committing to an affine probe \emph{collapses the irreducible coupling}, because
an affine field \textbf{factors through each vertex's own position}---reintroducing the
frozen, single-vertex reference that \S\,(why the quadric resists reduction) said mechanics
lacked. With $\mathbf{a} = (\mathrm{vec}(F),\mathbf{c})\in\mathbb{R}^{12}$, the displacement
at node $n$ is
\begin{equation}
  \mathbf{d}_n = F X_n + \mathbf{c}
    = \underbrace{[\,X_n^T\otimes I_3 \mid I_3\,]}_{P(X_n)\,\in\,\mathbb{R}^{3\times12}}
      \,\mathbf{a},
\end{equation}
depending only on node $n$'s rest position. The per-element probe energy is therefore
\begin{equation}
  W_e(\mathbf{a}) = \tfrac12\,\mathbf{a}^T G_e\,\mathbf{a},
  \
  G_e = V_e^T K_e V_e \in \mathbb{R}^{12\times12},
  \
  V_e = \begin{bmatrix} P(X_{n_1}) \\ P(X_{n_2}) \\ P(X_{n_3}) \end{bmatrix},
\end{equation}
and the global probe stiffness $G = \sum_e G_e$ \textbf{accumulates additively over
elements}---a small, fixed-size ($12\times12$, effective rank $6$), positive semi-definite
object that \emph{is} the homogenized stiffness in probe coordinates. The collapse cost is
the deviation of $G$ from its frozen fine-mesh target,
\begin{equation}
  E_{\text{collapse}}(\mathbf{x}^\star)
    = \big\| G_{\text{fine}} - G_{\text{after}}(\mathbf{x}^\star) \big\|_W^2
    = \Big\| \sum_{e\,\in\,\text{affected}}
        \big( G_e^{\text{old}} - G_e^{\text{new}}(\mathbf{x}^\star) \big) \Big\|_W^2,
\end{equation}
which is local, additive, fixed-size, and minimized over the merged position
$\mathbf{x}^\star$---structurally identical to a QEM scoring step. This is the Li-style
reduction: where Li attaches a fixed-size additive quadric $[\mathbf{p},\mathbf{u}]$ to each
vertex, here the per-vertex ``attribute'' is the vertex's probe-response signature, and the
metric is the energy metric $D$ rather than the identity.

\paragraph{The unifying picture.} The probe subspace dials the mechanical quadric between
two known forms:
\begin{itemize}
  \item \textbf{probe $=$ affine modes} $\rightarrow$ the signature is determined by
        position $\rightarrow$ the quadric reduces to (essentially) geometric QEM reweighted
        by the constitutive metric $D$; preserving homogenized stiffness under affine modes
        is mostly preserving area/shape, which geometric QEM already does;
  \item \textbf{probe $=$ low-frequency eigenmodes} $\rightarrow$ the signature is the
        mode-shape value per vertex $\rightarrow$ exactly Li's attribute quadric, but with
        the elastic mode shape as attribute and energy as metric---Nakagawa \& Kanai's
        ``modal strain as saliency'' turned into an actual quadric.
\end{itemize}
The mechanical $K_e$-in-the-$D$-metric is the energy-metric generalization that contains
both as probe-subspace special cases, with the Schur complement as the exact,
full-space limit.

\paragraph{Two honest caveats.} (i) $G_e(\mathbf{x}^\star)$ is not \emph{exactly} quadratic
in $\mathbf{x}^\star$, because moving the rest vertex changes $B$, the area, and hence $K_e$
nonlinearly; either freeze the fine reference and accept an approximate quadratic, or
evaluate at candidate placements (midpoint/endpoints) and take the minimum, as many QEM
variants do. (ii) The bargain is identical to geometric QEM's: response is preserved
\emph{within the probe subspace only}---off-probe (higher-frequency, non-affine) response is
discarded, just as geometric QEM preserves flatness but not higher-order curvature.

\subsection{The proposed pipeline}

We can now assemble the pieces into an end-to-end, QEM-shaped decimation pipeline. The
one-line summary: \emph{it is QEM, but the face quadric is replaced by the element stiffness
projected into a probe subspace ($G_e = V_e^T K_e V_e$), accumulated additively per vertex;
an edge collapse is priced by how much homogenized response on the probe it destroys
---equivalently, the energy in the differential mode it freezes out---and the merged vertex
is placed to minimize that loss.}

\paragraph{Phase 0 --- assemble the mechanical model.} On the fine mesh, assemble the two
element families into the global operator,
\begin{equation}
  K_f = \underbrace{\sum_{\text{tri}} A t\,B^T D B}_{\text{stretching (rank 3)}}
      \;+\; \underbrace{\sum_{\text{edge}} \tfrac12\,c\,\mathbf{g}\mathbf{g}^T}_{\text{bending hinge (rank 1)}}.
\end{equation}
This is the assembly view: both families write into the $1$-ring blocks; no
``triangle-vs-edge step'' decision is needed.

\paragraph{Phase 1 --- choose the probe subspace $V$.} This is where ``responsiveness'' is
\emph{defined}.
\begin{itemize}
  \item \textbf{Default:} the six affine / constant-strain modes $\rightarrow$ preserves the
        homogenized membrane stiffness (the patch test).
  \item \textbf{If bending matters:} enrich with the constant-curvature (quadratic) modes
        $\mathbf{u}(X)=\tfrac12 X^T H X + FX + \mathbf{c}$---affine probes are blind to bending.
  \item \textbf{If global compliance / known loads matter:} append low-frequency eigenmodes
        of $K_f$ or the load subspace.
\end{itemize}

\paragraph{Phase 2 --- build the mechanics-aware quadric (precompute).} Per element, project
the stiffness into probe coordinates, $G_e = V_e^T K_e V_e$, with $V_e$ sampling the probe
modes at the element's nodes through each node's rest position $P(X_n)$. Accumulate
additively to vertices, $G_v = \sum_{e \ni v} G_e$---the direct analog of QEM face-quadric
accumulation---and freeze the fine-mesh target $G_{\text{fine}}$.

\paragraph{Phase 3 --- greedy edge-collapse decimation.} Run the standard QEM heap loop. For
each candidate edge $(v_i, v_j)$:
\begin{enumerate}
  \item \textbf{Place} the merged vertex $\mathbf{x}^\star$ (and its probe signature) to
        minimize the mechanical cost---choosing it so the single vertex reproduces what the
        pair did on the probe, i.e.\ so the eliminated differential mode $\mathbf{d}_-$ stays
        small.
  \item \textbf{Score} with the composite cost
        \begin{equation}
          C \;=\; \underbrace{\big\|G_{\text{after}}(\mathbf{x}^\star) - G_{\text{fine}}\big\|_W^2}_{E_{\text{mech}}\ \text{(probe response lost)}}
            \;+\; \gamma\,E_{\text{triq}} \;+\; \alpha\,E_{\text{geom}},
        \end{equation}
        where $E_{\text{triq}}$ is the sliver penalty (required for a \emph{simulatable}
        output, per Nakagawa \& Kanai's ablation) and $E_{\text{geom}}$ is optional standard
        geometric QEM for visual fidelity.
  \item \textbf{Collapse} the heap minimum and \textbf{re-evaluate bending in the whole
        $1$-ring} (dihedral references move---bending is not collapse-invariant).
\end{enumerate}
Repeat to the target element budget. The merged vertex carries a position $\mathbf{x}^\star$
\emph{plus} a probe-response attribute---exactly Li's augmented vertex $[\mathbf{p},\mathbf{u}]$,
except the attribute is the elastic probe signature and the metric is $D$ rather than the
identity.

\paragraph{Forks still open.} The pipeline has a default at each branch, but the following
choices are genuine design/experiment forks:
\begin{center}
\renewcommand{\arraystretch}{1.3}
\begin{tabular}{|p{0.24\linewidth}|p{0.42\linewidth}|p{0.24\linewidth}|}
\hline
\textbf{Fork} & \textbf{Options} & \textbf{Lean} \\
\hline
cost backbone & portable additive $\|G_{\text{after}}-G_{\text{fine}}\|$ vs.\ exact
  differential-mode Schur per collapse & (a); validate against (b) \\
\hline
probe subspace & affine / $+$ curvature / $+$ eigenmodes / $+$ loads & affine first \\
\hline
bending bookkeeping & fold into assembled $G_e$ vs.\ distribute to hinge's $4$ verts & assembled \\
\hline
merged placement & closed-form min of $G$-deviation vs.\ candidate set (mid/endpoints) & candidate set \\
\hline
metric $W$ & identity / energy-weighted / mode-importance & identity \\
\hline
\end{tabular}
\end{center}
Two checks de-risk the whole idea: (1) on a flat patch, the affine-probe additive quadric
should reduce to roughly $D$-reweighted geometric QEM (sanity); (2) on a single $1$-ring,
the additive cost should track the exact differential-mode Schur cost across collapse
candidates (correctness).

\subsection{Dirty meshes and simplicial complexes}

Real assets---a dress with dangling straps, lace, beads, non-manifold seams, disconnected
trim---stress the pipeline at its foundation: assembly. The load-bearing fact is that
\textbf{finite-element assembly is purely combinatorial and additive}: $K = \sum_e K_e$
scatters each element's dense block into the global matrix indexed by that element's DOFs,
and never asks whether an edge is manifold or a surface is closed. Everything below is a
consequence of that, and of where the consequence frays.

\paragraph{Constructing $K$ on a dirty mesh.} Assembly always succeeds; the real content is
the rule mapping each simplex to an element family, plus a convention at non-manifold edges.
\begin{center}
\renewcommand{\arraystretch}{1.3}
\begin{tabular}{|p{0.30\linewidth}|p{0.30\linewidth}|p{0.30\linewidth}|}
\hline
\textbf{cell} & \textbf{element} & \textbf{quadric} \\
\hline
$0$-cell (bead) & point mass & --- (mass only) \\
\hline
$1$-cell (strap, lace) & bar / rod & rank-$1$ stretch ($+$ rod bending) \\
\hline
$2$-cell (panel) & membrane & rank-$3$ stretch; $+$ hinge on its edges \\
\hline
$3$-cell (solid) & P1 tetrahedron & rank-$6$ solid strain \\
\hline
\end{tabular}
\end{center}
Membrane stiffness assembles per triangle no matter how triangles are connected (a
non-manifold edge shared by $k$ panels just sums $k$ contributions into the shared DOFs). The
bending hinge is where non-manifoldness bites: a boundary edge ($1$ triangle) has \emph{no}
hinge (correct---a free hem has no bending partner), while a non-manifold edge ($k>2$
triangles) has an \textbf{ambiguous} dihedral angle and needs an explicit convention (skip /
sum all $\binom{k}{2}$ pairs / dominant pair). Dangling lower-dimensional simplices carry no
membrane and must be given bar/point elements if they are to carry load.

\paragraph{Is the operator ``proper''?} ``Proper'' splits into four senses that come apart on
dirty meshes.
\begin{itemize}
  \item \textbf{Symmetry $+$ PSD --- always holds.} Each $K_e$ is a $\tfrac12\|M\mathbf{d}\|^2$
        form, so $K=\sum K_e$ is symmetric PSD regardless of mesh quality. It is always a
        legitimate energy Hessian.
  \item \textbf{Null space --- grows, and the growth is diagnostic.} For a clean connected
        shell $\ker K$ is the six rigid modes. Each disconnected component adds six more;
        under-constraint adds \emph{spurious} zero modes (an isolated bead; a flat
        membrane patch, whose out-of-plane motion is zero-energy at linear order \emph{unless
        bending is present}; CST's missing drilling DOF; a bowtie vertex's counter-rotating
        cones). Hence bending is not optional for shells---it is what controls the
        membrane null space.
  \item \textbf{Conditioning --- dies on degeneracy.} Slivers give $K_e \sim 1/A \to \infty$:
        meaningless strain and a catastrophic condition number. This is why the triangle-quality
        term is load-bearing, not cosmetic.
  \item \textbf{Continuum consistency --- genuinely fails at non-manifold junctions.} There is
        no smooth $2$-manifold being discretized there, so ``convergence to the elasticity
        PDE'' is undefined. The operator remains a legitimate \emph{discrete} elastic energy,
        just not a Galerkin approximation of a shell.
\end{itemize}
The key reconciliation: our target $A^T K_c A = A^T K_f A$ references the fine
\emph{discrete} operator, so it stays well-defined even with no continuum limit. \textbf{The
mechanical-QEM objective degrades gracefully}---we never needed the continuum claim. The
fixes are exactly (i) a bending convention, (ii) null-space regularization, (iii) sliver
cleanup, and (iii) is already in the cost.

\paragraph{Collapse on simplicial complexes.} The three legs of the framework transfer
off-manifold. Assembly is combinatorial ($\checkmark$); the cost is linear algebra on $K$
($\checkmark$); and the affine probe still factors through vertex position at \emph{any}
cell's nodes, so $G_e = V_e^T K_e V_e$ accumulates additively across mixed-dimensional cells
---\textbf{the per-vertex additive quadric survives unchanged}. Mixed cells couple
automatically in assembly (a vertex shared by a panel and a strap sums both contributions),
which is just the ``fold bending and stretching into the assembled operator'' decision
generalized to bars and point masses. The topological scaffolding is prior art: Popovi\'c \&
Hoppe's \emph{Progressive Simplicial Complexes} (1997) generalizes vertex-split/edge-collapse
to arbitrary-dimension, arbitrary-topology complexes, and Garland \& Zhou's \emph{Quadric-Based
Simplification in Any Dimension} (2005) already puts the \emph{geometric} quadric on
$n$-simplices. The clean statement of our contribution: \textbf{inherit that contraction and
quadric machinery and swap the geometric quadric for the mechanical one.} Two caveats remain:
\begin{itemize}
  \item \textbf{Topological legality.} Simplicial contraction is governed by the
        \textbf{link condition}; violating it silently merges sheets, pinches, or kills
        handles. Controlled topology change (closing a sub-millimeter gap between panels) is
        often desirable, but must be a deliberate operator, not an accident.
  \item \textbf{Spurious null modes corrupt the cost.} A floppy/dangling region has a
        zero-energy mode, so collapsing it costs \emph{zero} mechanical error---``you cannot
        lose response you never had.'' The decimator will therefore aggressively delete exactly
        the dangling straps and flaps. This is why, on dirty assets, the geometric/feature
        term is \textbf{mandatory}: it protects mechanically-free but visually-important
        structure. The eigenmode probe inherits the same problem (the lowest modes of a dirty
        $K$ \emph{are} the floppy modes), which is a further reason to default to affine
        probes.
\end{itemize}

\paragraph{Volumetric decimation.} Extending from $2$-complexes (shells) to $3$-complexes
(solids) is, perhaps surprisingly, \emph{cleaner}. A P1 tetrahedron gives
$K_e = V\,B^T D B$ with the $6\times6$ solid constitutive $D$ and \textbf{no separate bending
family}---bending of a solid is just a strain distribution through its thickness, so the
awkward stretching/bending split vanishes: one element family, one quadric per tet. Solid
elasticity is a well-posed $3$D PDE with no non-manifold-junction ambiguity, so the patch test
and convergence are meaningful again, and interior-node elimination is textbook static
condensation (substructuring / superelements): \textbf{the exact Schur cost is on its firmest
footing volumetrically.} The operation gets stricter, though: tet edge-collapse must avoid
\textbf{inversion} (negative volume $\to$ indefinite $K_e$), so legality is the link condition
\emph{plus} a positive-volume constraint on all incident tets. And one usually wants to
preserve the boundary surface (visual) while coarsening the interior (mechanical)---a
two-criterion problem the simplicial-complex view handles natively: boundary triangles carry a
geometric quadric, interior tets carry the mechanical quadric, both accumulate to shared
vertices, the collapse cost sums them. The unifying structure is therefore a \textbf{single
abstract simplicial complex with a per-cell element type}, decimated by simplicial contraction
gated by (i) topological legality, (ii) geometric non-inversion, and (iii) the mechanical
cost. Surfaces and solids stop being separate pipelines.

\paragraph{The central hazard and a cheap diagnostic.} Across all of the above, the spurious
null space is the one thing we must \emph{design} rather than inherit. It is also cheaply
localizable: in $3$D,
\begin{equation}
  \big(\dim\ker K\big) - 6\,\big(\#\text{connected components}\big)
    \;=\; \#\,\text{spurious mechanism modes,}
\end{equation}
and visualizing \emph{where} those modes have support shows precisely the regions mechanics
would delete for free. Mitigations: a small regularization $K \to K + \epsilon M$ or a tiny
drilling/normal stiffness; always co-weighting a geometric/feature term; a cleanup pre-pass
(weld duplicates, drop zero-area/volume faces, snap near-coincident geometry) before any FEM.
The bottom line: the stiffness matrix can always be \emph{built} and is always a proper
\emph{discrete} energy; what a dirty mesh costs is continuum faithfulness (only at
non-manifold junctions, which our objective does not need) and conditioning (fixable by
cleanup).
